\section{Background and Related Work}
\label{sec:related}

\subsection{SORT, DeepSORT, and the Gating Contract}

Tracking-by-detection casts multi-object tracking as per-frame detection followed by frame-to-frame data association. SORT \cite{bewley2016sort} pairs a constant-velocity Kalman filter with greedy Hungarian assignment on intersection-over-union, prioritizing speed. DeepSORT \cite{wojke2017deepsort} augments this core with a learned appearance descriptor and a cosine-distance gate \cite{wojke2018cosine}, reducing identity switches under occlusion. Across both, the association core -- prediction, Hungarian matching, and the appearance gate -- is the component most often credited with tracking quality.

Less discussed, but central to this work, is that the original DeepSORT does not consume raw detections directly. Its public pipeline~\cite{wojke2017deepsortcode} applies an explicit \emph{gating contract} to the precomputed detection set before any of them reach the tracker: detections below a confidence threshold (\texttt{min\_confidence} $=0.3$) are discarded, boxes below a minimum height are removed, and per-frame non-maximum suppression (NMS) prunes overlapping responses. Only the survivors are wrapped as \texttt{Detection(tlwh, confidence, feature)} objects and passed to \texttt{predict()}/\texttt{update()}. This filter is a precision-oriented admission policy: it trades recall for a cleaner detection stream, suppressing the low-confidence and duplicate boxes that would otherwise instantiate spurious tracks. Modernization efforts that swap in a contemporary detector behind the same association core can silently \textbf{drop this gate}, since the new detector is typically queried in a streaming, ungated fashion; our audit shows this omission, rather than the association core, drives the most consequential behavior on dense clips.

\subsection{HOTA and the DetA/AssA Decomposition}

We evaluate with Higher Order Tracking Accuracy (HOTA) \cite{luiten2021hota} via TrackEval \cite{luiten2020trackeval}. HOTA's appeal here is structural: it factorizes geometrically into a detection term and an association term, $\mathrm{HOTA} \approx \sqrt{\mathrm{DetA} \cdot \mathrm{AssA}}$, where DetA scores how well predicted boxes localize ground-truth objects and AssA scores how consistently identities are maintained across the matched detections. This decomposition lets us attribute a change in tracking quality to detection versus association rather than reading a single conflated number, which is precisely the lever our study exercises: a modernization gain or a regression can be assigned to the DetA or the AssA factor, and a precision change in the detection stream registers directly in DetA. Earlier scalar metrics such as MOTA conflate these effects and weight identity errors lightly, obscuring the detection-driven regime we report. This detection-dominated regime is documented in the work that motivates HOTA. Luiten et al.~\cite{luiten2021hota} show that for the top MOT17 trackers detection error accounts for almost all of the variance in the legacy MOTA score (a linear fit gives $R^2 \approx 0.99$ against the detection term versus $R^2 \approx 0.24$ against the association term, with the detection-to-association error ratio ranging from roughly $42$ to $186$), the very imbalance HOTA was designed to correct. The dominance persists in modern systems: ByteTrack~\cite{zhang2022bytetrack} reaches strong accuracy with no appearance model at all, and StrongSORT~\cite{du2023strongsort} attributes more of its gains to a stronger detector than to its appearance and association refinements. An equal-magnitude change in the detection stream therefore moves the score more than a comparable association change -- the premise our decomposition (\cref{sec:decomp}) tests directly.

\subsection{ByteTrack and the Gating Mechanism}

ByteTrack \cite{zhang2022bytetrack} is the natural point of comparison, since it too revisits how low-confidence detections are handled, but its mechanism is the opposite of an admission gate and must be positioned carefully. ByteTrack \emph{keeps} low-confidence boxes rather than thresholding them away; crucially, it uses them only in a second association pass to recover \emph{existing} tracks, and never lets a low-confidence box \emph{instantiate} a new track (new tracks are seeded only from the remaining high-confidence detections, so ByteTrack already embodies a precision gate for instantiation). Our finding concerns ungated low-confidence detections that are free to spawn spurious new tracks, inflating identity counts and false positives on dense scenes. These are complementary mechanisms operating on the same raw material: ByteTrack constrains where weak detections may act (association only), whereas the dropped DeepSORT contract constrained whether they enter at all. ByteTrack is therefore not a counterexample to our result but an instance of the same underlying principle -- weak detections must be admitted under restriction -- realized at a different stage of the pipeline.

\subsection{Modern DeepSORT Descendants}

A family of trackers extends the DeepSORT design. StrongSORT \cite{du2023strongsort} upgrades the appearance model, motion compensation, and feature update; OC-SORT \cite{cao2023ocsort} reworks the motion model to be observation-centric and robust to nonlinear motion and occlusion; BoT-SORT \cite{aharon2022botsort} adds camera-motion compensation and a fused IoU--appearance cost. The lineage has continued through 2024--2025 along axes orthogonal to track admission: Deep OC-SORT \cite{maggiolino2023deepocsort} folds adaptive appearance and camera-motion compensation into the OC-SORT association cost, UCMCTrack \cite{yi2024ucmctrack} replaces IoU with a ground-plane Mahalanobis distance under uniform camera-motion compensation, and Hybrid-SORT \cite{yang2024hybridsort} fuses detector confidence, box height, and velocity direction as additional weak cues. These advances target the association core, and they generally inherit a confidence threshold as a fixed hyperparameter rather than treating it as the dominant, scene-dependent lever.

Crucially, the admission and track-initialization gate is not absent from these descendants; it is instantiated as an explicit, separately tuned hyperparameter. ByteTrack \cite{zhang2022bytetrack} seeds new tracks only from unmatched high-confidence detections, with the birth threshold set above its association threshold; BoT-SORT \cite{aharon2022botsort} exposes the same decision as a first-class \texttt{new\_track\_thresh}; OC-SORT \cite{cao2023ocsort} spawns tracklets only from detections passing a \texttt{det\_thresh}; and StrongSORT \cite{du2023strongsort} inherits DeepSORT's confidence gate (\texttt{min\_confidence}) and temporal confirmation gate (\texttt{n\_init}). Our contribution is therefore not the discovery that gating matters -- the field already treats the gate as a named, tuned parameter -- but the audit of a routine modernization in which this otherwise-explicit gate is silently dropped, together with a per-clip, precision-dependent quantification of what the omission costs. The per-clip precision-gating insight is in this sense orthogonal and composable: it does not compete with the association improvements above but specifies how the detection stream feeding any of them should be admitted, and is most decisive exactly where their motion and appearance machinery has the least to work with -- dense, low-resolution clips with unreliable detections.

A parallel line treats detection confidence as an actively engineered lever rather than a fixed threshold. BoostTrack \cite{stanojevic2024boosttrack} and BoostTrack++ \cite{stanojevic2025boosttrackpp} compute a detection--tracklet confidence score and use it to boost or reweight detections so that more genuine objects are admitted, while Localization-Guided Track \cite{meng2023lgtrack} decomposes detection confidence into classification and localization components to choose the association cost rather than discarding low-confidence boxes uniformly. That such effort is spent deliberately engineering confidence handling underscores how consequential the admission policy is, and how readily a naive modernization can mishandle it.

A complementary family, joint-detection-and-embedding trackers such as JDE \cite{wang2020jde} and FairMOT \cite{zhang2021fairmot}, learns detection and appearance together, arguing that anchor-based detection otherwise corrupts the re-identification features trained alongside it; our ground-truth-box isolation does not contradict this position, since it measures the descriptor's marginal value with detection \emph{held fixed} rather than claiming appearance is unimportant in general.

\subsection{ImageNet Features as a Re-identification Baseline}

Finally, generic ImageNet-pretrained features are a well-established re-identification baseline, routinely reported as the lower anchor that purpose-built descriptors such as OSNet \cite{zhou2019osnet} (here via torchreid \cite{zhou2019torchreid}) must clear. Our near-parity result between a non-REID-trained timm backbone \cite{wightman2019timm} and OSNet is therefore \emph{not} a claim that ImageNet features rival dedicated REID models in general. It is a controlled isolation: under ground-truth boxes the detection stage is held fixed, removing detector-induced variance so that association quality is measured directly, and in that regime the two appearance models are nearly indistinguishable on these short sequences. We report it to quantify how little of the modernization gain is attributable to the appearance descriptor once detection is controlled for, reinforcing the detection-dominated decomposition that motivates the rest of the paper.
